% Responsable: QUINTERO RICO LUIS DAVID.
\section{Resumen ejecutivo}
\label{sec:resumen-ejecutivo}

El Sistema de Gestión de Ayudantías de Cátedra (SGAC) organiza en una sola plataforma el proceso de ayudantías de la Universidad Técnica Estatal de Quevedo (UTEQ): convocatoria, postulación, verificación de requisitos, concurso de méritos y oposición, seguimiento de actividades y certificación. El análisis de requisitos previo describe un proceso apoyado en correos y archivos separados, en el que es difícil conocer el estado de una solicitud, conservar el historial de decisiones y consolidar informes~\cite{sgacAntecedente}. La entrevista de levantamiento del 21 de noviembre de 2026 añadió dos necesidades: que la consulta del estado sea ágil y que una caída del Sistema de Gestión Académica (SGA) no obligue al estudiante a postular de nuevo. La digitalización de la gestión universitaria exige integrar procesos, datos y responsables~\cite{fernandez2023digital,souli2025administracion}.

Usarán el sistema los estudiantes postulantes y ayudantes, los docentes, los coordinadores de carrera, la comisión de selección, el decano y el administrador. El estudiante podrá postular y seguir su trámite en línea; la coordinación y la comisión, verificar requisitos y registrar puntajes con una rúbrica común; y el docente, revisar las actividades y los informes del ayudante. Los doce requisitos funcionales (RF-01 a RF-12) se expresan como historias de usuario con criterios de aceptación verificables.

Se opta por un sistema distribuido por tres razones que fijan los requisitos no funcionales. La carga se concentra en los cinco días de inscripción: el escenario de diseño es de 350 postulaciones por convocatoria y 100 usuarios simultáneos en el pico (RNF-02), con una disponibilidad del 99\,\% en esa ventana (RNF-03). La verificación académica depende del SGA, que el equipo no controla: si no responde, la postulación debe conservarse en «Validación pendiente» (RNF-04) y el SGAC nunca escribe en él (RNF-10). Además, las funciones cambian a ritmos distintos y deben desplegarse por separado (RNF-11). ADR-001 adopta pocos servicios por capacidad de negocio (postulaciones, elegibilidad, evaluación, seguimiento y notificaciones) y ADR-002 define REST con JSON para las consultas y mensajería para los avisos.

El Proyecto Fin de Curso se desarrolla en cuatro entregas. E1 documenta el problema, los requisitos y la arquitectura, e incluye un esqueleto con los servicios de postulaciones y elegibilidad comunicados según un contrato versionado. E2 desarrollará los microservicios, sus API y contenedores; E3, los datos distribuidos y el procesamiento; y E4, la integración, las pruebas y la defensa del sistema. Durante E1 a E4 el SGA se representa con un servicio simulado de solo lectura. El asistente de inteligencia artificial y la firma electrónica quedan fuera del alcance.
